\scanpage{22}\mspage{21}
Si on pose $\deleted{\unsure{S|J}}F=$ frontière relative de $Y$ dans $X$, \deleted{\unsure{c.} F} Alors $F$ est fini, \inserted{modéré,} en pos. mod. p.r. à $Y$, \deleted{\unsure{donc|dans}} $F\subset Z$, \deleted{$\unsure{F\cap Z}$} \deleted{\unsure{est}} $F$ ne \deleted{$\unsure{Z}$} contient pas de pts isolé de $Z$.
\deleted{et on a}
\[
\deleted{\mathrm{Comp}}\ C_{\mathrm{mod}}(Y,X)=C_{\mathrm{mod}}(\deleted{\unsure{Z}}F,Z)
\]
\figure{\tikz{\draw (0,0) .. controls (.07,.02) and (.15,.10) .. (.21,.08) .. controls (.24,-.03) and (.38,-.01) .. (.55,.04) .. controls (.70,.11) and (.74,-.04) .. (.90,-.04) .. controls (1.22,-.05) and (1.33,.05) .. (1.48,.05) .. controls (1.66,-.04) and (1.78,-.08) .. (1.95,-.06) .. controls (2.16,-.05) and (2.27,.08) .. (2.41,.05) .. controls (2.56,-.05) and (2.81,-.05) .. (3.04,.00) .. controls (3.20,.02) and (3.32,.12) .. (3.42,.06) .. controls (3.50,-.08) and (3.80,-.02) .. (4.10,-.01);}}

(9 juin) Soit $\mathcal L$ un tronçon ordonné.

\deleted{Un sous-segment dans \uline{sous-tronçon} \uline{fermé} de $\mathcal L$ est un sous ens de la forme $[a,$}

\textcircled{1} Topologie sur $\mathcal L$ : base d'ouverts les $]a,b[$ (pour $a<b$), $\mathcal L_{<a}$\deleted{\unsure{,}}, $\mathcal L_{>b}$
[NB si $\mathcal L$ n'a pas de plus petit élément (resp. de plus grand élément) on peut omettre les $\mathcal L_{<a}$ (resp. les $\mathcal L_{>b}$). \inserted{Une extrémité de} \deleted{Pour que $\unsure{\{a\}}$ soit} $\mathcal L$ \unsure{—|(} (s'il y en a) est fermé, mais \deleted{fermé, il f et \unsure{suff} s. que}
On fera attention qu'en général, \deleted{les $\unsure{\{a\}}$} \deleted{\unsure{si $a$ est|si $a$ et}} pour $a\in\mathcal L$ $\{a\}$ \uline{n'est pas} \inserted{\uline{néc$^{\mathrm t}$}} \uline{fermé}. Si $\mathcal L$ est l'ens des lieux d'une \deleted{\unsure{segm}} ficelle topologique$^{\ast}$, et alors $\deleted{\unsure{a}}b\in\overline{\{\deleted{\unsure{b}}a\}}$ i.e. $\overline{\{b\}}\subset\overline{\{a\}}$, équivaut à $b\supset a$ (en tt que parties de \unsure{$\mathscr H$|$\mathcal H$|$H$})

% pale vertical watermark
\rotatebox{90}{\textsc{\unsure{Strong Extra}}}
